\newpage\thispagestyle{empty}
\chapter*{Abstract}
\addcontentsline{toc}{chapter}{Abstract}
\vspace{0.5cm}

Central Bank Digital Currencies (CBDCs) promise to modernise
payment infrastructure while extending financial inclusion to
unbanked populations. However, the transition from physical cash
to digital tokens introduces a critical privacy challenge: every
transaction leaves a metadata trail that can expose citizens to
re-identification even when no names are recorded. This risk is
particularly significant in Nepal, where the Nepal Rastra Bank
(NRB) is actively exploring a domestic token-based CBDC for
deployment across diverse geographic and demographic contexts.

Existing approaches to CBDC privacy comparison are qualitative in
nature, offering no formal metric by which policymakers can
measure, compare, or regulate metadata de-anonymization risk
across different cryptographic scheme choices. This thesis
addresses that gap by developing a formal, normalised privacy risk
model grounded in information-theoretic principles.

The model introduces three core components: a single-attribute
uniqueness score $s(a_i)$, a pairwise interaction coefficient
$\lambda_{ij}$ that captures amplification effects between
attribute combinations, and a normalised risk score
$R_{\mathrm{norm}} = R_{\mathrm{raw}} / R_{\mathrm{max}}$,
where $R_{\mathrm{max}}(m) = m + \binom{m}{2}$ is derived from
first principles as the proven theoretical upper bound. The model
is applied to six cryptographic schemes --- Plain Hash Token,
Blind Signature, Hybrid Scheme, Ring Signature, Stealth Address,
and ZKP Token --- across seven metadata attributes, including two
Nepal-specific fields: payment channel and transaction zone.

Results demonstrate that the six schemes form a clean,
non-overlapping risk ladder from $R_{\mathrm{norm}} = 0.317$
(Plain Hash Token) to $R_{\mathrm{norm}} = 0.035$ (ZKP Token).
Population scale provides no meaningful privacy improvement for
weak schemes, with Plain Hash Token at $N = 100{,}000$ remaining
$5.5\times$ riskier than ZKP Token. Nepal's demographic clustering
provides passive privacy protection at domestic scale. Statistical
validation across 30 random seeds confirms the scheme ranking is
robust. A micro-scale analysis at $N = 50$ reveals critical risk
at village pilot scale, with all non-ZKP schemes exceeding
$R_{\mathrm{norm}} = 0.48$.

The findings provide NRB with a formal, quantitative basis for
cryptographic scheme selection prior to national CBDC deployment,
with Ring Signature combined with a full mitigation stack
identified as the recommended minimum for resource-constrained
deployment contexts.